% Ad Familiares 13.29 --- headnote
% ad-familiares-13-29, 47 BC, Rome

\textit{Cicero to Lucius Munatius Plancus, written from
Rome late in 47 BC or at the beginning of 46 BC (the
manuscript dateline: \textit{Scr. Romae vel ex. a. 707
(47) vel in. a. 708 (46)}). Plancus, son of Cicero's
old friend Lucius Munatius Plancus the elder, was at
this date a partisan of Caesar's --- he would shortly be
governor of one of the Gauls --- though the suppleness
of allegiance that defines his later career is already
visible in the way Cicero approaches him here: as a
younger man with influence at Caesar's ear, and one
whose patronage Cicero is prepared to draw on.}

\textit{This is no formulaic commendaticia but a sustained
letter of recommendation, framed by an unusually long
personal preamble about the inherited tie between the
two families and their shared literary interests. The
beneficiary is Gaius Ateius Capito, a close friend who
stood by Cicero through the political reversals of the
50s and the civil war, and what Cicero is asking for is
substantive: that Plancus, using Caesar's favour,
secure for Capito the inheritance of his kinsman Titus
Antistius. The body of the letter is a careful
exoneration of Antistius's conduct in Macedonia ---
caught by Pompey's arrival as quaestor, he had no real
choice; he kept aloof from the camp, hid himself in the
interior, was treated mildly by Caesar after Pharsalus,
and died of illness on his way back to Rome. Sect.~7
contains Cicero's own discreet but pointed
self-defence about his conduct during the war: that he
behaved more moderately on Pompey's side than anyone
else thanks largely to Capito's counsel, a remark that
also serves the recommendation. The closing tier-marker
--- ``in such terms as I could press no other claim
with greater care or greater earnestness'' --- is among
the warmest formulae of the genre.}
